% Prior-corrected CAWPE for HIVE-COTE 2.0.
%
% Self-contained section, intended to be \input into a paper. Requires:
%   \usepackage{amsmath, amssymb, booktabs, graphicx}
% and a \newtheorem{proposition}{Proposition} declaration if one is not already present.

\section{Prior-Corrected Combination}
\label{sec:prior-correction}

\subsection{Background and notation}

HIVE-COTE 2.0 (HC2)~\cite{middlehurst2021hc2} combines $M=4$ component classifiers, each
built on a different representation of the series: the Shapelet Transform Classifier
(STC), the Diverse Representation Canonical Interval Forest (DrCIF), Arsenal, and the
Temporal Dictionary Ensemble (TDE). Let $\mathcal{D}=\{(\mathbf{x}_i,y_i)\}_{i=1}^{n}$ be
a training set of series $\mathbf{x}_i$ with labels $y_i \in \mathcal{C} =
\{1,\dots,C\}$, and let $\hat{p}_j(c \mid \mathbf{x})$ denote the probability estimate
that component $j$ assigns to class $c$.

Components are combined by CAWPE~\cite{large2019cawpe}, which weights each component by
its training accuracy estimate raised to a fixed exponent $\alpha$ (default $\alpha=4$).
Writing $\hat{a}_j$ for the out-of-bag accuracy estimate of component $j$, the weights
and the combined posterior are
\begin{equation}
  w_j = \hat{a}_j^{\,\alpha},
  \qquad
  \hat{p}(c \mid \mathbf{x})
  = \frac{\sum_{j=1}^{M} w_j \, \hat{p}_j(c \mid \mathbf{x})}
         {\sum_{c' \in \mathcal{C}} \sum_{j=1}^{M} w_j \, \hat{p}_j(c' \mid \mathbf{x})}.
  \label{eq:cawpe}
\end{equation}
The prediction is then $\hat{y}(\mathbf{x}) = \arg\max_{c} \hat{p}(c \mid \mathbf{x})$,
with ties broken uniformly at random.

\subsection{The decision rule implied by a loss}

The rule $\hat{y} = \arg\max_c \hat{p}(c \mid \mathbf{x})$ minimises the expected
$0/1$ loss, and therefore targets accuracy. It is not the optimal rule for balanced
accuracy. Let $\pi_c = \Pr(y=c)$ denote the class prior, estimated from the training set
as $\hat{\pi}_c = n_c / n$ with $n_c = |\{i : y_i = c\}|$. Balanced accuracy is the mean
of the per-class recalls,
\begin{equation}
  \mathrm{BA} = \frac{1}{C} \sum_{c \in \mathcal{C}}
    \Pr\!\left(\hat{y}(\mathbf{x}) = c \mid y = c\right),
  \label{eq:ba}
\end{equation}
which is exactly the accuracy that would be obtained if every class were equally likely.
Under a uniform prior $\pi'_c = 1/C$, Bayes' rule gives
\begin{equation}
  \Pr\!\left(y = c \mid \mathbf{x};\, \pi'\right)
  \;\propto\; p(\mathbf{x} \mid y = c) \, \pi'_c
  \;\propto\; \frac{\hat{p}(c \mid \mathbf{x})}{\pi_c},
  \label{eq:bayes}
\end{equation}
since $p(\mathbf{x}\mid y=c) \propto \hat{p}(c\mid\mathbf{x})/\pi_c$. Dividing the
combined posterior by the training prior therefore converts the accuracy-optimal rule
into the balanced-accuracy-optimal one. This is a change of decision rule only: the
probability estimates themselves, and every component, are left untouched.

\subsection{Tempered prior correction}

Equation~\eqref{eq:bayes} assumes $\hat{p}(c\mid\mathbf{x})$ is calibrated, so that the
only discrepancy between the two rules is the prior. In practice the CAWPE posterior of
Eq.~\eqref{eq:cawpe} is a renormalised weighted vote rather than a calibrated
probability, and full correction over-shoots. We therefore introduce a single
temperature parameter $\beta \in [0,1]$ and define the \emph{prior-corrected} posterior
\begin{equation}
  \boxed{\;
  \hat{p}^{\,\beta}(c \mid \mathbf{x})
  = \frac{\hat{p}(c \mid \mathbf{x}) \; \hat{\pi}_c^{\,-\beta}}
         {\sum_{c' \in \mathcal{C}} \hat{p}(c' \mid \mathbf{x}) \; \hat{\pi}_{c'}^{\,-\beta}},
  \qquad
  \hat{y}^{\,\beta}(\mathbf{x}) = \arg\max_{c \in \mathcal{C}} \hat{p}^{\,\beta}(c \mid \mathbf{x}).
  \;}
  \label{eq:tempered}
\end{equation}
Equivalently, in the log domain the correction is an additive class-wise offset,
$\log \hat{p}^{\,\beta}(c\mid\mathbf{x}) = \log \hat{p}(c\mid\mathbf{x}) - \beta \log
\hat{\pi}_c + \mathrm{const}$, i.e.\ a shift of the decision boundary against frequent
classes in proportion to their log-frequency. The parameter $\beta$ interpolates between
the two rules: $\beta=0$ recovers standard HC2 exactly, and $\beta=1$ gives the Bayes
rule for a uniform prior. Three properties follow directly.

\begin{proposition}[Invariance on balanced data]
If $\hat{\pi}_c = 1/C$ for all $c$, then $\hat{p}^{\,\beta} = \hat{p}$ for every $\beta$,
and the predictions are unchanged.
\end{proposition}
\noindent The factor $\hat{\pi}_c^{-\beta} = C^{\beta}$ is constant in $c$ and cancels in
the normalisation. The correction is thus inactive by construction on balanced problems
and cannot harm them.

\begin{proposition}[Invariance under degenerate estimates]
\label{prop:degenerate}
If $\hat{p}(\cdot\mid\mathbf{x})$ is a one-hot vector, then $\hat{y}^{\,\beta}(\mathbf{x})
= \hat{y}(\mathbf{x})$ for every $\beta$.
\end{proposition}
\noindent Scaling the single non-zero entry by a positive constant cannot change which
entry is largest. Prior correction therefore applies only to classifiers that emit
genuinely soft estimates, and has no effect on hard-voting classifiers. We return to
this in Section~\ref{sec:pc-comparison}.

\begin{proposition}[Cost]
Computing $\hat{y}^{\,\beta}$ for all test cases requires $O(nC)$ additional operations
and no refitting, no access to the training series, and no additional model storage.
\end{proposition}

\subsection{Experimental setup}

We evaluate on the 112 univariate UCR datasets~\cite{dau2019ucr} over 30 stratified
resamples, reporting the mean over resamples per dataset and comparing classifiers with
the Wilcoxon signed-rank test over the 112 paired dataset means. All HC2 variants are
constructed from the same stored component predictions, so differences are attributable
to the combination scheme alone and not to component variance. Datasets are grouped by
imbalance ratio $\rho = \max_c n_c / \min_c n_c$: balanced ($\rho = 1$, 27 datasets),
mild ($1 < \rho < 2$, 52), and imbalanced ($\rho \geq 2$, 33, of which 11 have
$\rho \geq 5$).

\subsection{Results}

Table~\ref{tab:beta} reports both metrics as a function of $\beta$. The key observation
is that mild correction is \emph{not} a trade between the two metrics: for
$\beta \leq 0.75$ both accuracy and balanced accuracy improve simultaneously, and only
full correction ($\beta=1$) sacrifices accuracy. This indicates that the CAWPE posterior
is genuinely miscalibrated towards frequent classes, rather than merely being tuned to a
different operating point.

\begin{table}[t]
\centering
\caption{Effect of the prior-correction temperature $\beta$ on HC2. Mean over 112 UCR
datasets and 30 resamples; $\Delta$ in percentage points against $\beta=0$; $p$ from a
Wilcoxon signed-rank test over the 112 paired dataset means.}
\label{tab:beta}
\begin{tabular}{lcccccc}
\toprule
& \multicolumn{3}{c}{Accuracy} & \multicolumn{3}{c}{Balanced accuracy} \\
\cmidrule(lr){2-4}\cmidrule(lr){5-7}
$\beta$ & \% & $\Delta$ & $p$ & \% & $\Delta$ & $p$ \\
\midrule
$0$ (HC2)      & 89.301 & --- & --- & 87.244 & --- & --- \\
$0.25$         & 89.429 & $+0.129$ & $0.0001$ & 87.640 & $+0.396$ & $<0.0001$ \\
$\mathbf{0.50}$& \textbf{89.469} & $\mathbf{+0.169}$ & $\mathbf{0.0017}$ & 87.982 & $+0.738$ & $<0.0001$ \\
$0.75$         & 89.341 & $+0.041$ & $0.0446$ & 88.227 & $+0.982$ & $<0.0001$ \\
$1.00$         & 88.991 & $-0.310$ & $0.1808$ & \textbf{88.380} & $\mathbf{+1.135}$ & $<0.0001$ \\
\bottomrule
\end{tabular}
\end{table}

The effect is concentrated where the theory predicts. Table~\ref{tab:groups} breaks the
$\beta=0.5$ result down by imbalance. Balanced datasets are unchanged to three decimal
places, as required by Proposition~1. Gains in balanced accuracy grow monotonically with
skew, reaching $+3.309$ percentage points on the most imbalanced group. Accuracy gains
peak on mildly imbalanced data and turn slightly negative at $\rho \geq 5$, where a fixed
$\beta = 0.5$ over-corrects.

\begin{table}[t]
\centering
\caption{Change from prior correction at $\beta=0.5$, by imbalance ratio $\rho$.
Percentage points, mean over 30 resamples.}
\label{tab:groups}
\begin{tabular}{lccc}
\toprule
Group & $n$ & $\Delta$ Accuracy & $\Delta$ Balanced accuracy \\
\midrule
Balanced ($\rho=1$)      & 27 & $+0.000$ & $+0.000$ \\
Mild ($1<\rho<2$)        & 52 & $+0.258$ & $+0.402$ \\
Imbalanced ($\rho\geq2$) & 33 & $+0.166$ & $+1.871$ \\
Severe ($\rho\geq5$)     & 11 & $-0.149$ & $+3.309$ \\
\bottomrule
\end{tabular}
\end{table}

\subsubsection{A constant outperforms a tuned parameter}

Since $\beta$ is a free parameter, an obvious refinement is to select it per dataset. We
compared a fixed $\beta = 0.5$ against selecting $\beta$ from
$\{0, 0.125, \dots, 1\}$ on the combined \emph{training} posterior, built from the
components' own out-of-bag estimates, so that no test data is used. Tuning is worse than
the constant on both metrics (Table~\ref{tab:tuning}): paired against the fixed setting
it loses $0.046$ percentage points of accuracy and wins on only 25 of 112 datasets
($p = 0.059$). The cause is a systematic bias in the selection signal: training-based
selection chooses a mean $\beta$ of $0.167$ with a \emph{median} of $0.019$, that is, it
usually declines to correct at all, against an oracle mean of $0.223$. The training
posterior is evaluated against the same class distribution the components were fitted to,
so the uncorrected rule appears better on training data than it is on test data.

\begin{table}[t]
\centering
\caption{Selecting $\beta$ per dataset against a fixed value. Oracle selects $\beta$
per dataset on test accuracy and is an upper bound, not a method.}
\label{tab:tuning}
\begin{tabular}{lcccc}
\toprule
& \multicolumn{2}{c}{Accuracy} & \multicolumn{2}{c}{Balanced accuracy} \\
\cmidrule(lr){2-3}\cmidrule(lr){4-5}
Selection of $\beta$ & \% & $\Delta$ & \% & $\Delta$ \\
\midrule
None ($\beta=0$)          & 89.301 & --- & 87.244 & --- \\
Fixed $\beta=0.5$         & \textbf{89.469} & $\mathbf{+0.169}$ & 87.982 & $+0.738$ \\
Tuned on train accuracy   & 89.423 & $+0.122$ & 87.664 & $+0.419$ \\
Tuned on train bal.\ acc. & 89.155 & $-0.146$ & \textbf{88.166} & $\mathbf{+0.921}$ \\
\midrule
Oracle (test accuracy)    & 89.831 & $+0.530$ & 88.148 & $+0.903$ \\
\bottomrule
\end{tabular}
\end{table}

The oracle row shows that roughly three times the gain remains available to a perfect
selection rule. Closing that gap requires a selection signal other than the training
posterior; we leave this open.

\subsubsection{Comparison with reweighting the ensemble}
\label{sec:pc-comparison}

Prior correction alters the decision rule; the obvious alternative is to alter the
weights $w_j$. We find the latter to be structurally limited. At $\alpha=4$ over training
accuracies that lie in a narrow band, CAWPE assigns weight shares between $22.6\%$ and
$28.9\%$, against $25\%$ for uniform weighting, so the weights have little dynamic range
to exploit. Accordingly, replacing $\hat{a}_j$ with the true test accuracy of each
component -- an oracle -- improves balanced accuracy by only $0.05$ percentage points.
Searching the full weight simplex per dataset, with weights selected on one stratified
half of the test set and scored on the other, gains $0.33$ percentage points of balanced
accuracy and is not significant (better on 51 of 112 datasets, $p = 0.211$). Prior
correction at $\beta = 0.5$ exceeds this ceiling on both metrics while requiring no
tuning data.

Finally, we compare against MultiRocket-Hydra~\cite{dempster2023hydra}, which is not
significantly different from HC2 on this archive. Holding the components fixed, prior
correction changes the balanced-accuracy comparison from not significant to significant:
uncorrected HC2 is better on 56 of 112 datasets ($+0.665$ percentage points, $p=0.298$),
while at $\beta=0.5$ it is better on 72 of 112 ($+1.403$ points, $p<0.0001$). Both rows
use identical component predictions and differ only in the decision rule, so the change
is attributable to the correction.

We make no corresponding claim for accuracy from this comparison. In our reconstruction
HC2 is already significantly more accurate than MultiRocket-Hydra before any correction
($+0.903$ points, 64 of 112, $p=0.0119$), so the correction does not confer that
significance.

Using the published reference results as the canonical build, neither comparison is
significant before correction ($+0.742$ points on accuracy, $p=0.101$; $+0.474$ points
on balanced accuracy, $p=0.782$). Transferring the per-dataset effect of $\beta=0.5$
additively onto those results estimates $+0.911$ points on accuracy ($p=0.010$) and
$+1.211$ points on balanced accuracy ($p=0.0002$). We report this as an estimate rather
than a measurement: published reference results provide aggregate accuracies only, with
no probability estimates to which the correction could be applied.

We note that the comparison is fair in a specific and asymmetric sense:
applying the same correction to MultiRocket-Hydra changes nothing whatsoever, on either
metric and at every $\beta$, because its ridge-based head emits one-hot estimates for
$100\%$ of test cases and Proposition~\ref{prop:degenerate} applies. The ability to
benefit from post-hoc calibration is a property of producing soft estimates, which HC2's
weighted probabilistic combination provides and a hard-voting pipeline does not. This is
a property of the implementations compared rather than a statement that convolutional
pipelines cannot be calibrated.

\subsection{Threats to validity}
\label{sec:pc-threats}

Component predictions are stochastic, and independently rebuilt components do not
reproduce a stored run exactly. Taking the published reference results as the canonical
build, our reconstruction of uncorrected HC2 is significantly \emph{more} accurate than
published HC2 ($+0.159$ percentage points, better on 64 of 112 datasets, $p=0.0015$),
because it inherits component runs postdating the reference HC2 run. The discrepancy is
comparable in size to the effect we report. Comparisons between $\beta$ settings are
paired on identical component predictions and are unaffected, but absolute comparisons
against an external classifier are: uncorrected HC2 against MultiRocket-Hydra on accuracy
gives $p=0.101$ using published results and $p=0.0119$ using our reconstruction. Our
MultiRocket-Hydra results agree with the published reference to within $0.002$ percentage
points, so the discrepancy is specific to the HC2 build.

Consequently, the claims in Section~\ref{sec:pc-comparison} that rest on an external
comparison should be read as indicative. The directly measured, build-independent result
is the paired improvement of Table~\ref{tab:beta}. A definitive external comparison
requires rebuilding all four components and the corrected ensemble within a single run,
which we leave to future work.

The value $\beta = 0.5$ was selected on the same 112 datasets on which it is evaluated,
so the reported accuracy gain of $0.169$ percentage points should be treated as mildly
optimistic; validation on an independent collection is required before the magnitude is
relied upon. The direction of the effect, and the invariance results of Propositions~1
and~\ref{prop:degenerate}, do not depend on this choice. Separately, the correction may
be treating a symptom: three of HC2's four components emit near-degenerate estimates
(Arsenal is one-hot on $80\%$ of cases, and TDE accumulates member weights into the
winning class only, so it has no probability estimates at all), which plausibly explains
why the combined posterior is sharper and more majority-biased than a calibrated one.
Calibrating the components before combination may subsume the correction proposed here.
